\documentclass[11pt,a4paper]{article}

% ═══════════════════════════════════════════════════════════════════════════════
% PACKAGES
% ═══════════════════════════════════════════════════════════════════════════════
\usepackage[utf8]{inputenc}
\usepackage[T1]{fontenc}
\usepackage[french]{babel}
\usepackage{amsmath,amssymb}
\usepackage{siunitx}
\usepackage{physics}
\usepackage{tikz}
\usepackage{pgfplots}
\usepackage{xcolor}
\usepackage{tcolorbox}
\usepackage{booktabs}
\usepackage{geometry}
\usepackage{fancyhdr}
\usepackage{hyperref}
\usepackage{pifont}

% Configuration
\geometry{margin=2.5cm}
\pgfplotsset{compat=1.18}
\usetikzlibrary{arrows.meta, positioning, shapes.geometric, patterns, decorations.pathmorphing, calc, backgrounds}

% Couleurs personnalisées
\definecolor{colEmis}{HTML}{3498DB}
\definecolor{colEntres}{HTML}{2ECC71}
\definecolor{colAbsorbes}{HTML}{E74C3C}
\definecolor{colCompton}{HTML}{E67E22}
\definecolor{colTheo}{HTML}{F39C12}
\definecolor{colEau}{HTML}{85C1E9}
\definecolor{colAir}{HTML}{F9E79F}

% En-tête
\pagestyle{fancy}
\fancyhf{}
\rhead{Simulation Am-241 -- Geant4}
\lhead{Problème de pollution Compton}
\rfoot{Page \thepage}

% ═══════════════════════════════════════════════════════════════════════════════
\begin{document}

\begin{center}
    {\LARGE\bfseries Explication du problème de pollution Compton}\\[0.3cm]
    {\large dans les spectres d'absorption Am-241}\\[0.5cm]
    {\normalsize Simulation Geant4 -- Analyse des résultats}
\end{center}

\vspace{0.5cm}

% ═══════════════════════════════════════════════════════════════════════════════
\section{Contexte et observation du problème}
% ═══════════════════════════════════════════════════════════════════════════════

Dans la simulation Geant4 de l'irradiation par une source Am-241, on observe une anomalie dans le spectre des gammas entrant dans l'eau : certaines raies (notamment 43.4 et 55.6~keV) semblent absentes ou très atténuées, alors que des counts supplémentaires apparaissent dans des bins adjacents.

\begin{tcolorbox}[colback=red!5!white, colframe=red!75!black, title=\textbf{Observation problématique}]
\begin{itemize}
    \item Bin 54-56~keV : 21 gammas émis $\rightarrow$ \textbf{146 gammas entrés} (+125 !)
    \item Bin 56-58~keV : 0 gamma émis $\rightarrow$ \textbf{200 gammas entrés} (+200 !)
    \item Bin 58-60~keV : 35\,794 gammas émis $\rightarrow$ 29\,182 gammas entrés (-6\,612)
\end{itemize}
\end{tcolorbox}

% ═══════════════════════════════════════════════════════════════════════════════
\section{Géométrie de la simulation}
% ═══════════════════════════════════════════════════════════════════════════════

\begin{figure}[h!]
\centering
\begin{tikzpicture}[scale=1.0, >=Stealth]
    
    % Fond
    \fill[colAir!30] (-1, -0.5) rectangle (6, 4);
    \fill[colEau!50] (6, -0.5) rectangle (9.5, 4);
    
    % Source Am-241
    \node[circle, fill=yellow!80!orange, minimum size=1cm, draw=orange!80!black, thick] 
        (source) at (0, 1.75) {\small Am-241};
    
    % Étiquettes des zones
    \node[above] at (2.5, 4) {\textbf{Air (25 mm)}};
    \node[above] at (7.75, 4) {\textbf{Eau (3 mm)}};
    
    % Lignes de séparation
    \draw[very thick, dashed] (6, -0.5) -- (6, 4);
    
    % Dimensions
    \draw[<->, thick] (0, -0.3) -- (6, -0.3) node[midway, below] {25 mm};
    \draw[<->, thick] (6, -0.3) -- (9, -0.3) node[midway, below] {3 mm};
    
    % Gamma direct (59.5 keV)
    \draw[->, colEmis, very thick] (0.5, 2.5) -- (8.5, 2.5);
    \node[colEmis, right] at (8.6, 2.5) {\small 59.5 keV};
    \fill[colEmis] (0.5, 2.5) circle (3pt);
    
    % Gamma diffusé Compton
    \draw[->, colEmis, thick] (0.5, 1.5) -- (3.5, 1.5);
    \fill[colEmis] (0.5, 1.5) circle (3pt);
    
    % Point de diffusion Compton
    \fill[colCompton] (3.5, 1.5) circle (5pt);
    \node[colCompton, above right] at (3.5, 1.7) {\small Compton};
    
    % Gamma après diffusion
    \draw[->, colCompton, very thick, dashed] (3.5, 1.5) -- (8.5, 0.8);
    \node[colCompton, right] at (8.6, 0.8) {\small 55 keV};
    
    % Électron Compton
    \draw[->, gray, thick] (3.5, 1.5) -- (4.5, 2.8);
    \node[gray, above] at (4.5, 2.8) {\small $e^-$};
    
    % Angle de diffusion
    \draw[thick] (4.2, 1.5) arc (0:-25:0.7);
    \node at (4.8, 1.2) {\small $\theta$};
    
    % Légende
    \node[anchor=north west] at (-0.8, 0) {
        \begin{tabular}{cl}
            \tikz\draw[->, colEmis, very thick] (0,0) -- (0.5,0); & Gamma direct \\
            \tikz\draw[->, colCompton, very thick, dashed] (0,0) -- (0.5,0); & Gamma diffusé \\
        \end{tabular}
    };
    
\end{tikzpicture}
\caption{Géométrie de la simulation : les gammas traversent 25~mm d'air avant d'atteindre l'eau. Certains subissent une diffusion Compton dans l'air et arrivent avec une énergie réduite.}
\label{fig:geometrie}
\end{figure}

% ═══════════════════════════════════════════════════════════════════════════════
\section{Physique de la diffusion Compton}
% ═══════════════════════════════════════════════════════════════════════════════

La diffusion Compton est l'interaction d'un photon avec un électron quasi-libre. Le photon est dévié d'un angle $\theta$ et perd une partie de son énergie.

\subsection{Formule de l'énergie après diffusion}

L'énergie du photon après diffusion Compton est donnée par :

\begin{equation}
\boxed{E' = \frac{E_0}{1 + \dfrac{E_0}{m_e c^2}(1 - \cos\theta)}}
\label{eq:compton}
\end{equation}

où :
\begin{itemize}
    \item $E_0$ = énergie initiale du photon (59.5~keV pour la raie principale Am-241)
    \item $E'$ = énergie après diffusion
    \item $m_e c^2$ = 511~keV (masse de l'électron)
    \item $\theta$ = angle de diffusion
\end{itemize}

\subsection{Application numérique pour la raie 59.5~keV}

\begin{figure}[h!]
\centering
\begin{tikzpicture}
\begin{axis}[
    width=12cm, height=7cm,
    xlabel={Angle de diffusion $\theta$ (°)},
    ylabel={Énergie après Compton (keV)},
    xmin=0, xmax=180,
    ymin=40, ymax=62,
    grid=both,
    grid style={gray!30},
    legend pos=south west,
    legend style={fill=white, fill opacity=0.8},
]

% Courbe Compton
\addplot[colEmis, very thick, domain=0:180, samples=100] 
    {59.5 / (1 + (59.5/511)*(1 - cos(x)))};
\addlegendentry{$E'(\theta)$ pour $E_0 = 59.5$~keV}

% Lignes horizontales pour les raies
\addplot[colEntres, thick, dashed] coordinates {(0, 59.5) (180, 59.5)};
\addlegendentry{Raie 59.5~keV}

\addplot[colCompton, thick, dashed] coordinates {(0, 55.6) (180, 55.6)};
\addlegendentry{Raie 55.6~keV}

\addplot[colAbsorbes, thick, dashed] coordinates {(0, 43.4) (180, 43.4)};
\addlegendentry{Raie 43.4~keV}

% Zone de pollution (48-58 keV)
\addplot[colCompton, fill opacity=0.2] fill between[of=A and B] 
    coordinates {(0,48) (180,48) (180,58) (0,58)};

% Points remarquables
\addplot[only marks, mark=*, mark size=3pt, black] 
    coordinates {(45, 57.5) (60, 56.2) (90, 53.3) (120, 50.7)};

% Annotations
\node[anchor=west] at (axis cs:47, 57.5) {\small $\theta=45°$};
\node[anchor=west] at (axis cs:62, 56.2) {\small $\theta=60°$};
\node[anchor=west] at (axis cs:92, 53.3) {\small $\theta=90°$};
\node[anchor=west] at (axis cs:122, 50.7) {\small $\theta=120°$};

% Zone pollution label
\node[colCompton, fill=white, fill opacity=0.8] at (axis cs:140, 53) {\small Zone pollution};

\end{axis}
\end{tikzpicture}
\caption{Énergie d'un gamma de 59.5~keV après diffusion Compton en fonction de l'angle $\theta$. La zone orangée (48-58~keV) correspond aux bins qui seront "pollués" par les gammas diffusés.}
\label{fig:compton_curve}
\end{figure}

\begin{table}[h!]
\centering
\caption{Énergie après diffusion Compton pour $E_0 = 59.5$~keV}
\label{tab:compton}
\begin{tabular}{ccc}
\toprule
\textbf{Angle $\theta$} & \textbf{Énergie $E'$} & \textbf{Bin affecté} \\
\midrule
0° & 59.5~keV & 58-60~keV (raie 59.5) \\
30° & 58.6~keV & 58-60~keV \\
45° & 57.5~keV & 56-58~keV \\
\rowcolor{colCompton!20} 60° & 56.2~keV & 56-58~keV $\leftarrow$ \textit{pollution} \\
\rowcolor{colCompton!20} 90° & 53.3~keV & 52-54~keV $\leftarrow$ \textit{pollution} \\
\rowcolor{colCompton!20} 120° & 50.7~keV & 50-52~keV $\leftarrow$ \textit{pollution} \\
180° & 48.3~keV & 48-50~keV \\
\bottomrule
\end{tabular}
\end{table}

% ═══════════════════════════════════════════════════════════════════════════════
\section{Illustration du problème}
% ═══════════════════════════════════════════════════════════════════════════════

\begin{figure}[h!]
\centering
\begin{tikzpicture}[scale=0.9]

% ─────────────────────────────────────────────────────────────────────────────
% Panel gauche : Spectre ÉMIS
% ─────────────────────────────────────────────────────────────────────────────
\begin{scope}[shift={(0,0)}]
    \draw[->] (0,0) -- (7,0) node[right] {E (keV)};
    \draw[->] (0,0) -- (0,5) node[above] {Counts};
    
    % Titre
    \node[anchor=south, font=\bfseries] at (3.5, 5.2) {Spectre ÉMIS};
    
    % Axe X
    \foreach \x/\label in {1/44, 2/48, 3/52, 4/56, 5/60, 6/64} {
        \draw (\x, 0) -- (\x, -0.1) node[below, font=\small] {\label};
    }
    
    % Barres du spectre émis
    \fill[colEmis, opacity=0.8] (0.7, 0) rectangle (1.3, 0.3);  % 43.4 keV (faible)
    \fill[colEmis, opacity=0.8] (3.7, 0) rectangle (4.3, 0.1);  % 55.6 keV (très faible)
    \fill[colEmis, opacity=0.8] (4.7, 0) rectangle (5.3, 4.5);  % 59.5 keV (dominant)
    
    % Annotations
    \node[colEmis, font=\small] at (1, 0.6) {43.4};
    \node[colEmis, font=\small] at (4, 0.4) {55.6};
    \node[colEmis, font=\small] at (5, 4.8) {59.5};
    
    % Légende
    \node[anchor=north west, font=\small] at (0.2, 4.8) {Raies Am-241};
\end{scope}

% Flèche
\draw[->, very thick, gray] (7.5, 2.5) -- (9, 2.5) node[midway, above] {\small Air 25mm};

% ─────────────────────────────────────────────────────────────────────────────
% Panel droit : Spectre ENTRÉ (pollué)
% ─────────────────────────────────────────────────────────────────────────────
\begin{scope}[shift={(9.5,0)}]
    \draw[->] (0,0) -- (7,0) node[right] {E (keV)};
    \draw[->] (0,0) -- (0,5) node[above] {Counts};
    
    % Titre
    \node[anchor=south, font=\bfseries, colAbsorbes] at (3.5, 5.2) {Spectre ENTRÉ (pollué)};
    
    % Axe X
    \foreach \x/\label in {1/44, 2/48, 3/52, 4/56, 5/60, 6/64} {
        \draw (\x, 0) -- (\x, -0.1) node[below, font=\small] {\label};
    }
    
    % Zone de pollution Compton
    \fill[colCompton, opacity=0.2] (1.5, 0) rectangle (4.5, 3);
    \node[colCompton, font=\small\itshape] at (3, 2.5) {Pollution};
    \node[colCompton, font=\small\itshape] at (3, 2.0) {Compton};
    
    % Barres du spectre entré
    \fill[colAbsorbes, opacity=0.8] (0.7, 0) rectangle (1.3, 0.25);  % 43.4 keV
    \fill[colCompton, opacity=0.8] (1.7, 0) rectangle (2.3, 0.4);   % 48-50 keV (Compton)
    \fill[colCompton, opacity=0.8] (2.7, 0) rectangle (3.3, 0.9);   % 52-54 keV (Compton)
    \fill[colCompton, opacity=0.8] (3.7, 0) rectangle (4.3, 1.5);   % 56-58 keV (Compton + 55.6)
    \fill[colAbsorbes, opacity=0.8] (4.7, 0) rectangle (5.3, 3.8);  % 59.5 keV (réduit)
    
    % Annotations
    \draw[->, thick, colCompton] (3, 3.3) -- (2.5, 1.2);
    \draw[->, thick, colCompton] (3, 3.3) -- (3.5, 1.7);
    \node[colCompton, font=\small, above] at (3, 3.3) {Gammas 59.5 keV};
    \node[colCompton, font=\small, above] at (3, 3.0) {diffusés};
    
    % Différence
    \draw[<->, thick, red] (5.5, 4.5) -- (5.5, 3.8);
    \node[red, font=\small, right] at (5.5, 4.15) {-15\%};
\end{scope}

\end{tikzpicture}
\caption{Comparaison schématique des spectres émis et entré dans l'eau. Les gammas de 59.5~keV qui diffusent Compton dans l'air arrivent avec une énergie réduite et "polluent" les bins 48-58~keV.}
\label{fig:pollution}
\end{figure}

% ═══════════════════════════════════════════════════════════════════════════════
\section{Conséquence sur l'analyse}
% ═══════════════════════════════════════════════════════════════════════════════

\subsection{Le problème : identification par énergie mesurée}

Lorsqu'on construit un histogramme en utilisant l'énergie du gamma \textbf{au moment où il entre dans l'eau}, on commet une erreur d'identification :

\begin{figure}[h!]
\centering
\begin{tikzpicture}[
    box/.style={rectangle, draw, thick, minimum width=3cm, minimum height=1cm, align=center},
    arrow/.style={->, thick, >=Stealth}
]

% Gamma émis à 59.5 keV
\node[box, fill=colEmis!20] (emis) at (0, 0) {Gamma émis\\$E_0 = 59.5$~keV};

% Diffusion Compton
\node[circle, fill=colCompton!50, minimum size=1.5cm] (compton) at (4, 0) {Compton\\$\theta=60°$};

% Gamma avec énergie réduite
\node[box, fill=colCompton!20] (diffuse) at (8, 0) {Gamma diffusé\\$E' = 56$~keV};

% Entrée dans l'eau
\node[box, fill=colEau!50] (eau) at (12, 0) {Entre dans l'eau\\mesure : 56~keV};

% Flèches
\draw[arrow] (emis) -- (compton);
\draw[arrow] (compton) -- (diffuse);
\draw[arrow] (diffuse) -- (eau);

% Deux méthodes d'identification
\node[box, fill=red!20, text width=4cm] (methode1) at (6, -3) 
    {\textbf{Histogramme en E}\\Compté dans bin\\54-56~keV\\$\rightarrow$ \textcolor{red}{FAUX !}};

\node[box, fill=green!20, text width=4cm] (methode2) at (12, -3) 
    {\textbf{Statistiques par raie}\\Identifié comme\\raie 59.5~keV\\$\rightarrow$ \textcolor{green!50!black}{CORRECT}};

% Flèches vers méthodes
\draw[arrow, red, dashed] (eau.south) -- (methode1.north);
\draw[arrow, green!50!black] (eau.south) -- (methode2.north);

% Labels
\node[font=\small\itshape, gray] at (2, 0.8) {dans l'air};

\end{tikzpicture}
\caption{Un gamma émis à 59.5~keV mais diffusé Compton est mal identifié si on utilise son énergie mesurée. La méthode correcte utilise l'énergie \textbf{initiale} (LineID dans gamma\_lines).}
\label{fig:methodes}
\end{figure}

\subsection{La solution : identification par énergie initiale}

Le TTree \texttt{gamma\_lines} dans le fichier ROOT identifie chaque gamma par son énergie \textbf{initiale} (au moment de l'émission), pas son énergie au moment de l'entrée dans l'eau.

\begin{tcolorbox}[colback=green!5!white, colframe=green!75!black, title=\textbf{Solution recommandée}]
\textbf{Utiliser les statistiques par raie} (TTree \texttt{gamma\_lines}) :
\begin{itemize}
    \item Chaque gamma est identifié par son \texttt{LineID} (0-11)
    \item L'identification se fait par l'énergie \textbf{initiale}
    \item Même si le gamma diffuse Compton, il reste identifié correctement
    \item $\Rightarrow$ Pas de pollution entre les raies !
\end{itemize}
\end{tcolorbox}

% ═══════════════════════════════════════════════════════════════════════════════
\section{Comparaison des deux méthodes de visualisation}
% ═══════════════════════════════════════════════════════════════════════════════

\begin{figure}[h!]
\centering
\begin{tikzpicture}[scale=0.85]

% ─────────────────────────────────────────────────────────────────────────────
% Méthode INCORRECTE : Histogramme en énergie
% ─────────────────────────────────────────────────────────────────────────────
\begin{scope}[shift={(0,0)}]
    \node[anchor=south, font=\large\bfseries, red!70!black] at (3.5, 5.5) 
        {$\times$ INCORRECT : Histogramme en énergie};
    
    \draw[->] (0,0) -- (7.5,0) node[right] {E (keV)};
    \draw[->] (0,0) -- (0,5) node[above] {Counts};
    
    % Axe X
    \foreach \x/\label in {0.5/10, 1.5/20, 2.5/30, 3.5/40, 4.5/50, 5.5/60, 6.5/100} {
        \draw (\x, 0) -- (\x, -0.1) node[below, font=\tiny] {\label};
    }
    
    % Barres (spectre entré - pollué)
    \fill[colEmis, opacity=0.8] (0.3, 0) rectangle (0.7, 1.5);
    \fill[colEmis, opacity=0.8] (0.8, 0) rectangle (1.2, 3.5);
    \fill[colEmis, opacity=0.8] (1.3, 0) rectangle (1.7, 4.2);
    \fill[colEmis, opacity=0.8] (1.8, 0) rectangle (2.2, 2.0);
    \fill[colEmis, opacity=0.8] (2.3, 0) rectangle (2.7, 1.2);
    \fill[colCompton, opacity=0.6] (3.3, 0) rectangle (3.7, 0.3);  % 43 keV + Compton
    \fill[colCompton, opacity=0.6] (4.0, 0) rectangle (4.4, 0.5);  % Compton
    \fill[colCompton, opacity=0.6] (4.5, 0) rectangle (4.9, 0.8);  % Compton
    \fill[colCompton, opacity=0.6] (5.0, 0) rectangle (5.4, 1.2);  % 55.6 + Compton
    \fill[colEmis, opacity=0.8] (5.5, 0) rectangle (5.9, 4.5);     % 59.5 keV
    
    % Zone pollution
    \draw[colCompton, thick, dashed] (3.2, 0) rectangle (5.5, 1.5);
    \node[colCompton, font=\small] at (4.3, 1.8) {Pollution Compton};
    
    % Problème
    \node[red, font=\small, align=center] at (4.3, 3.5) {Les raies 43 et 56~keV\\sont noyées !};
\end{scope}

% ─────────────────────────────────────────────────────────────────────────────
% Méthode CORRECTE : Barres par raie
% ─────────────────────────────────────────────────────────────────────────────
\begin{scope}[shift={(9,0)}]
    \node[anchor=south, font=\large\bfseries, green!50!black] at (3.5, 5.5) 
        {$\checkmark$ CORRECT : Barres par raie};
    
    \draw[->] (0,0) -- (7.5,0) node[right] {Raie};
    \draw[->] (0,0) -- (0,5) node[above] {Counts};
    
    % Barres par raie (données correctes)
    \fill[colEntres, opacity=0.8] (0.5, 0) rectangle (1.0, 1.5);
    \fill[colEntres, opacity=0.8] (1.2, 0) rectangle (1.7, 3.5);
    \fill[colEntres, opacity=0.8] (1.9, 0) rectangle (2.4, 4.2);
    \fill[colEntres, opacity=0.8] (2.6, 0) rectangle (3.1, 2.0);
    \fill[colEntres, opacity=0.8] (3.3, 0) rectangle (3.8, 1.2);
    \fill[colEntres, opacity=0.8] (4.0, 0) rectangle (4.5, 0.15);  % 33 keV
    \fill[colEntres, opacity=0.8] (4.7, 0) rectangle (5.2, 0.08);  % 43 keV
    \fill[colEntres, opacity=0.8] (5.4, 0) rectangle (5.9, 0.03);  % 56 keV
    \fill[colEntres, opacity=0.8] (6.1, 0) rectangle (6.6, 4.5);   % 59.5 keV
    
    % Labels raies
    \node[font=\tiny, rotate=45, anchor=east] at (0.75, -0.3) {X$_{Ll}$};
    \node[font=\tiny, rotate=45, anchor=east] at (1.45, -0.3) {X$_{L\alpha}$};
    \node[font=\tiny, rotate=45, anchor=east] at (2.15, -0.3) {X$_{L\beta}$};
    \node[font=\tiny, rotate=45, anchor=east] at (2.85, -0.3) {X$_{L\gamma}$};
    \node[font=\tiny, rotate=45, anchor=east] at (3.55, -0.3) {$\gamma$26};
    \node[font=\tiny, rotate=45, anchor=east] at (4.25, -0.3) {$\gamma$33};
    \node[font=\tiny, rotate=45, anchor=east] at (4.95, -0.3) {$\gamma$43};
    \node[font=\tiny, rotate=45, anchor=east] at (5.65, -0.3) {$\gamma$56};
    \node[font=\tiny, rotate=45, anchor=east] at (6.35, -0.3) {$\gamma$59};
    
    % Avantage
    \node[green!50!black, font=\small, align=center] at (3.5, 3.5) 
        {Chaque raie est\\comptée séparément\\$\Rightarrow$ Pas de pollution};
\end{scope}

\end{tikzpicture}
\caption{Comparaison des deux méthodes de visualisation. À gauche : l'histogramme en énergie est pollué par les Compton. À droite : les barres par raie (depuis \texttt{gamma\_lines}) donnent les bonnes statistiques.}
\label{fig:comparaison}
\end{figure}

% ═══════════════════════════════════════════════════════════════════════════════
\section{Résumé et recommandations}
% ═══════════════════════════════════════════════════════════════════════════════

\begin{tcolorbox}[colback=blue!5!white, colframe=blue!75!black, title=\textbf{Résumé}]
\begin{enumerate}
    \item Les gammas de 59.5~keV traversent 25~mm d'air avant d'atteindre l'eau
    \item Certains subissent une \textbf{diffusion Compton} dans l'air
    \item Après diffusion : $E' = \frac{59.5}{1 + \frac{59.5}{511}(1-\cos\theta)}$ keV
    \item Ces gammas arrivent dans l'eau avec une énergie \textbf{réduite} (48-58~keV)
    \item Sur l'histogramme en énergie, ils "polluent" les bins des raies faibles
\end{enumerate}
\end{tcolorbox}

\vspace{0.5cm}

\begin{tcolorbox}[colback=green!5!white, colframe=green!75!black, title=\textbf{Recommandation}]
Pour l'analyse des taux d'absorption par raie :
\begin{itemize}
    \item[\ding{51}] Utiliser le TTree \texttt{gamma\_lines} (statistiques par raie)
    \item[\ding{51}] Faire des graphiques en \textbf{barres par raie}, pas des histogrammes en énergie
    \item[\ding{55}] Ne PAS utiliser \texttt{hGammaEnteringWater} pour l'analyse par raie
\end{itemize}
\end{tcolorbox}

\end{document}
